% This template is designed to offer an aesthetically pleasing resume that adheres to a formal and institutional tone, making it suitable for applications to companies and research centers requiring a high degree of professionalism. Navy blue has been chosen as the primary color to align with these objectives.
% Author: Christian Maria Giannetti

%----------------------------------------------------------------------------------------
%   Packages And Other Document Configurations
%----------------------------------------------------------------------------------------

\documentclass{resume} % Use the custom resume.cls style

% Document margins
\usepackage[left=0.75in,top=0.6in,right=0.75in,bottom=0.6in]{geometry}

% Color and hyperlink packages
\usepackage{xcolor}
\usepackage{hyperref}

% Footnote and margin adjustment packages
\usepackage{footnote}
\usepackage{changepage}

% Fontawesome package for icons
\usepackage{fontawesome}

% Tabularx package for custom tables
\usepackage{tabularx}

% Define navyblue color
\definecolor{navyblue}{RGB}{0,54,123}

%----------------------------------------------------------------------------------------
%   Customizations
%----------------------------------------------------------------------------------------

% Define italicitem, bolditem, and plainitem commands
\newcommand{\italicitem}[1]{\item{\textit{#1}}}
\newcommand{\bolditem}[1]{\item{\textbf{#1}}}
\newcommand{\plainitem}[1]{\item{#1}}

% Define user-friendly link command for hyperlinks
\newcommand{\link}[2]{{\href{#1}{#2}}}

\newcommand{\entry}[2]{#1 & #2 \tabularnewline} % Defines an entry with two arguments: #1 for the first column and #2 for the second column

%----------------------------------------------------------------------------------------
%   Define envsection command for defining a new environment section
%----------------------------------------------------------------------------------------

\newcommand{\tableEnv}[2]{%
  \begin{rSection}{#1} % Begin rSection with the given name
    \begin{adjustwidth}{0.0in}{0.1in} % Set the left and right margins
      \begin{tabularx}{\linewidth}{@{} >{\bfseries}l @{\hspace{6ex}} X @{}}
        #2 % Print the content inside the tabularx environment
      \end{tabularx}
    \end{adjustwidth}
  \end{rSection}
}

%----------------------------------------------------------------------------------------
%   Begin document
%----------------------------------------------------------------------------------------

% Set name with navyblue color
\name{\color{navyblue} Giulio Tognetto}

\begin{document}

\printPersonalInfo{
  \personalInfo{\tag{Residence/domicile}\info{Campiglia dei Berici (VI), Italy}}
  \personalInfo{\tag{E-mail}\info{contact@giuliotognetto.dev} \infoSeparator\tag{Telephone number}\info{+39 347 347 7917}}
  \personalInfo{\tag{GitHub}\info{\link{https://github.com/GiulioTognetto}{github.com/GiulioTognetto}}}
}

%----------------------------------------------------------------------------------------
%   Profile section
%----------------------------------------------------------------------------------------

\begin{rSection}{About me}
    \begin{adjustwidth}{0.0in}{0.1in}
        Persona creativa e curiosa, orientata allo sviluppo software full-stack e all'innovazione tecnologica. Appassionato di programmazione e di soluzioni originali, unisco competenze tecniche e approccio pratico per realizzare progetti concreti.
    \end{adjustwidth}
\end{rSection}

%----------------------------------------------------------------------------------------
%   Education section
%----------------------------------------------------------------------------------------

\begin{rSection}{Education}

    % Bachelor's degree entry
    \begin{rSubsectionNoBullet}{\bf Laurea Triennale in Scienze Informatiche}{Università degli Studi di Padova (UniPD)}{1° anno (In corso)}{2026 - Presente}
        \italicitem{Curriculum in Scienze Informatiche}
    \end{rSubsectionNoBullet}

    % High School diploma entry
    \begin{rSubsectionNoBullet}{\bf Diploma ITI - Informatica e Telecomunicazioni}{IIS Euganeo - Este (PD)}{Quotazione: 86/100}{2021 - 2026}
        \italicitem{Specializzazione in Informatica e Telecomunicazioni}
    \end{rSubsectionNoBullet}

\end{rSection}

%----------------------------------------------------------------------------------------
%   Work experience section
%----------------------------------------------------------------------------------------

\begin{rSection}{Work experience}

    % Freelance experience entry
    \begin{rSubsection}{Sviluppatore Freelance / Progetti Indipendenti}{2025 - Presente}{Sviluppo Web e Automazione}{Remoto}
        \item Sviluppo di soluzioni software custom, script e applicazioni web per clienti internazionali e nazionali tramite piattaforme di freelancing (Fiverr, Addlance).
        \item Gestione autonoma del ciclo di vita del progetto, dal requisito del cliente al deployment e consegna.
    \end{rSubsection}

    % First work experience entry
    \begin{rSubsection}{Dev Code S.r.l.}{07/2025 - 21/07/2025}{PCTO: Programmazione Python per test del software}{Este (PD), Italy}
        \item Sviluppo di script in Python per l'automazione dei test del software aziendale.
        \item Collaborazione attiva con il team di sviluppo, applicando capacità di analisi, ascolto e lavoro di gruppo.
    \end{rSubsection}

\end{rSection}

%----------------------------------------------------------------------------------------
% Technical skills section
%----------------------------------------------------------------------------------------

\tableEnv{Technical skills}{
    \entry{Programming Languages}{Python, JavaScript, TypeScript, Java, C, C++, Rust, SQL (SQLite, PostgreSQL, MariaDB), Bash, \LaTeX, HTML, CSS, Markdown}
    \entry{Frameworks \& Libraries}{Vue.js, Nuxt.js, Svelte, Astro, Flask.py, FastAPI.py, Express.js, Tauri.rs, eGUI.rs, PyTorch, TensorFlow, Selenium}
    \entry{Tools \& Platforms}{VS Code, Qt Designer, Unity, Godot, Blender, Git, GitHub Desktop, NPM, Cloudflare (CDN, DB, serverless Workers)}
}

%----------------------------------------------------------------------------------------
% Language proficiencies section
%----------------------------------------------------------------------------------------

\tableEnv{Language proficiencies}{
    \entry{Italian}{Madrelingua}
    \entry{English}{Livello B1 (Certificazione Cambridge English)}
    \entry{French}{Livello A2 (scolastico)}
}

%----------------------------------------------------------------------------------------
% Soft skills section
%----------------------------------------------------------------------------------------

\tableEnv{Soft skills}{
    \entry{Competenze organizzative}{Ascolto attivo, comunicazione chiara, organizzazione del lavoro su più task e problem solving analitico applicato all'automazione.}
}

\end{document}